% Esquisse d'un programme (cote 119, copy 2): archive cote 119, site page 103 (PDF page 104), typescript page 25.
% typescript. Draft compared: esquisse/pages/esq_119_027.tex (copy 1, site 27).
% AI-transcribed by Claude Opus 5.5 (claude-opus-5-5, adaptive thinking, effort high), Message Batches API, batch esq119_full (msgbatch_015xqDHkTumDysRF1qngZhRJ), 2026-11-13; not yet checked by eye.
\scanpage{119/103}
\mspage{25}

temps me semble mûr pour réécrire une version nouvelle, en style
moderne, du classique livre de Klein sur l'icosaèdre et les autres
polyèdres pythagoriciens. Ecrire un tel exposé sur les 2-polyèdres
réguliers serait une magnifique occasion pour un jeune chercheur
de se familiariser aussi bien avec la géométrie des polyèdres et
leurs liens avec les géométries sphèrique, euclidienne, hyperbo-%
lique, et avec les courbes algébriques, qu'avec le langage et les
techniques de base de la géométrie algébrique moderne. S'en trou-%
vera-t-il un un jour pour saisir cette occasion ?

5. \qquad Je voudrais maintenant dire quelques mots sur certaines
réflexions qui m'ont fait comprendre le besoin de fondements nou-%
veaux pour la topologie ``géométrique'', dans une direction toute
différente de la notion de topos, et indépendante même des besoins
de la géométrie algébrique dite ``abstraite'' (sur des corps et an-%
neaux de base généraux). Le problème de départ, qui a commencé à
m'intriguer il doit y avoir une quinzaine d'années déjà, était
celui de définir une théorie de ``dévissage'' des structures stra-%
tifiées, pour les reconstituer, par un procédé canonique, à partir
de ``pièces de construction'' canoniquement déduits de la structure
donnée. Probablement l'exemple principal qui m'avait alors amené
à cette question était celui de la stratification canonique d'une
variété algèbrique singulière (ou d'un espace analytique complexe
ou réel singulier) par la suite décroissante de ses ``lieux singuliers''
successifs. Mais je devais sans doute pressentir déjà l'ubiquité
des structures stratifiées dans pratiquement tous les domaines de
la géométrie (que d'autres sûrement ont vu clairement bien avant
moi). Depuis, j'ai vu apparaître de telles structures, notamment,
dans toute situation de ``modules'' pour des objets géométriques sus-%
ceptibles non seulement de variation continue, mais en même temps
de phénomènes de ``degénérescence'' (ou de ``spécialisation'') - les
strates correspondant alors aux divers ``niveaux de singularité''
(ou aux types combinatoires associés) pour les objets considérés.

\margin{103}
